\documentclass[11pt,a4paper]{article}

% ---------------------------------------------------------
% PACKAGES
% ---------------------------------------------------------
\usepackage[utf8]{inputenc}
\usepackage[T1]{fontenc}
\usepackage{lmodern}
\usepackage{amsmath, amssymb, amsfonts}
\usepackage{bm}
\usepackage{graphicx}
\usepackage{hyperref}
\usepackage{authblk}
\usepackage{geometry}
\usepackage{cite}
\usepackage{physics}
\usepackage{tensor}
\usepackage{enumitem}

\geometry{margin=1in}

% ---------------------------------------------------------
% TITLE
% ---------------------------------------------------------
\title{\textbf{Companion XIV: Primordial Black Holes as CPT-Sector Objects in the Relaxation-Driven Cyclic Cosmology (RDCC)}}

\author[1]{Michael Lehmann}
\affil[1]{\small Independent Researcher, Relaxation-Driven Cyclic Cosmology (RDCC) Framework\\
\texttt{mi.lehmann@gmx.de}}

\date{\small \today}

% ---------------------------------------------------------
% DOCUMENT START
% ---------------------------------------------------------
\begin{document}

\maketitle

\begin{abstract}
Primordial black holes (PBHs) provide a uniquely sensitive probe of early-Universe physics, constraining the thermal history, the growth of perturbations, and the collapse thresholds during radiation domination. In the Relaxation-Driven Cyclic Cosmology (RDCC), the bipartite CPT-symmetric global state, the sectoral temperature asymmetry, and the universal relaxation mechanism jointly reshape the conditions for PBH formation. In this Companion we develop a sectoral interpretation of PBHs within RDCC: PBHs arise predominantly in the CPT-conjugate (mirror) sector, where the lower temperature, modified relaxation history, and altered perturbation dynamics naturally enhance PBH formation. The visible sector, by contrast, exhibits a strong suppression of PBHs due to relaxation-induced damping and an elevated effective collapse threshold. We derive the CPT-sector PBH formation mechanism, compute the resulting mass function, and identify the observational signatures associated with this sectoral asymmetry. This Companion establishes PBHs as a natural and robust component of the CPT sector and provides a new observational window into the bipartite structure and relaxation-driven dynamics of RDCC.
\end{abstract}

% ========================================================
% 1 INTRODUCTION
% ========================================================

\section{Introduction}

Primordial black holes (PBHs) provide a uniquely sensitive probe of the earliest stages of cosmic
evolution. Their formation depends on the thermal history, the growth of primordial
perturbations, and the collapse thresholds during radiation domination. In standard cosmology,
PBHs require either enhanced primordial fluctuations, non-Gaussian features, or exotic early-time
dynamics to form in appreciable abundance.

In the Relaxation-Driven Cyclic Cosmology (RDCC), the situation is fundamentally different.
The bipartite CPT-symmetric global state, the sectoral temperature asymmetry, and the
relaxation-induced modification of perturbation dynamics jointly reshape the conditions for PBH
formation. The relaxation mechanism suppresses the growth of super-horizon perturbations in
the visible sector, raising the effective collapse threshold and strongly inhibiting PBH formation.
At the same time, the CPT-conjugate sector evolves with a lower temperature, reduced
relaxation-induced damping, and a larger horizon mass at a given scale factor.

These structural differences lead to a robust sectoral asymmetry: PBHs form naturally and
efficiently in the CPT-conjugate sector, while their formation in the visible sector is exponentially
suppressed. This interpretation arises directly from the bipartite global state and the universal
relaxation mechanism introduced in Companion~X, and it is consistent with the sectoral
thermodynamics developed in Companion~XIII.

The goal of this Companion is to develop the PBH sector of RDCC in a unified and
model-independent manner. Section~2 introduces PBHs as CPT-sector objects and explains the
suppression mechanism in the visible sector. Section~3 derives the PBH formation mechanism in
the CPT sector. Section~4 develops the modified collapse threshold and the resulting PBH mass
function. Section~5 analyzes the thermodynamic and dynamical properties of CPT-sector PBHs.
Section~6 presents the observational signatures, and Section~7 provides forecasts for future
surveys.

Together, these results establish PBHs as a natural and robust component of the CPT sector
and provide a new observational window into the bipartite structure and relaxation-driven
dynamics of RDCC.

% ========================================================
% 2 PBHs AS CPT-SECTOR OBJECTS
% ========================================================

\section{PBHs as CPT-Sector Objects}

The Relaxation-Driven Cyclic Cosmology (RDCC) is built on a bipartite, CPT-symmetric global
state in which two thermodynamically distinct sectors evolve on the same spacetime geometry.
A visible-sector observer corresponds to a projection onto the sector with a forward
thermodynamic arrow, while the CPT-conjugate sector evolves with decreasing entropy. The two
sectors interact only through gravity and the relaxation dynamics.

A key structural result of Companion~X is the persistent temperature asymmetry between the
sectors,
\begin{equation}
\frac{T_{B}}{T_{A}} \simeq 0.5,
\end{equation}
which remains stable throughout the early Universe. This temperature asymmetry has direct
implications for PBH formation, as both the collapse threshold and the horizon mass depend
sensitively on the thermal history.

PBH formation requires large density fluctuations to exceed a critical collapse threshold at
horizon re-entry. In the visible sector, the relaxation-induced damping of super-horizon
perturbations increases the effective collapse threshold, suppressing PBH formation. In the
CPT-conjugate sector, the lower temperature and reduced relaxation-induced damping lower the
collapse threshold and increase the horizon mass, enhancing PBH formation.

These structural differences lead to three robust conclusions:

\begin{enumerate}[label=(\alph*)]
\item \textbf{PBHs are strongly suppressed in the visible sector.}  
Relaxation-induced damping raises the effective collapse threshold, preventing overdensities
from reaching the required amplitude for collapse.

\item \textbf{PBHs form naturally in the CPT-conjugate sector.}  
The colder temperature and reduced damping lower the collapse threshold and increase the
horizon mass, making PBH formation efficient.

\item \textbf{PBHs are inherently CPT-sector objects in RDCC.}  
Their formation mechanism, abundance, and mass spectrum follow directly from the
thermodynamic and relaxation properties of the CPT sector.
\end{enumerate}

This interpretation resolves several conceptual tensions. It explains the absence of PBHs in the
visible sector without fine-tuning, provides a natural origin for PBHs in the RDCC framework,
and connects PBH physics directly to the bipartite structure of the global state.

In the next section, we derive the CPT-sector PBH formation mechanism and quantify the
degree of suppression in the visible sector.

% ========================================================
% 3 CPT-SECTOR PBH FORMATION MECHANISM
% ========================================================

\section{CPT-Sector PBH Formation Mechanism}

Primordial black holes (PBHs) form when large density fluctuations re-enter the horizon during
radiation domination and exceed a critical collapse threshold. In standard cosmology, this
threshold is typically of order
\begin{equation}
\delta_{\mathrm{std}} \sim 0.4 - 0.5,
\end{equation}
and PBH formation requires either enhanced primordial fluctuations or non-Gaussian features.

In the Relaxation-Driven Cyclic Cosmology (RDCC), the early-Universe dynamics are modified
by the universal relaxation mechanism introduced in Companion~X. The relaxation rate
\(\Gamma(a) = \alpha(a) H(a)\) alters the growth of super-horizon perturbations and increases the
effective collapse threshold in the visible sector. As a result, PBH formation in the visible sector
is strongly suppressed.

In the CPT-conjugate sector, the situation is fundamentally different. The lower temperature,
reduced relaxation-induced damping, and altered perturbation dynamics naturally enhance PBH
formation. Three structural ingredients govern the PBH formation mechanism in the CPT sector:

\begin{enumerate}[label=(\alph*)]

\item \textbf{Reduced relaxation-induced damping.}  
The relaxation coefficient satisfies
\begin{equation}
\alpha_{B}(a) < \alpha_{A}(a),
\end{equation}
due to the lower temperature of sector~B. This reduces the damping of super-horizon
perturbations and lowers the effective collapse threshold.

\item \textbf{Larger horizon mass.}  
The colder temperature implies a slower expansion rate and therefore a larger horizon mass
at a given scale factor:
\begin{equation}
M_{H}^{(B)}(a) \simeq 4\, M_{H}^{(A)}(a).
\end{equation}
This shifts the PBH mass spectrum in sector~B toward higher masses.

\item \textbf{Lower collapse threshold.}  
The reduced pressure in the colder sector lowers the collapse threshold:
\begin{equation}
\delta_{c}^{(B)} < \delta_{c}^{(A)}.
\end{equation}
This exponentially enhances PBH formation in sector~B.
\end{enumerate}

Combining these effects yields a natural and robust PBH formation mechanism in the CPT
sector. PBHs arise as a direct consequence of the bipartite global state and the relaxation
dynamics, without requiring enhanced primordial fluctuations or exotic early-Universe physics.

In contrast, the visible sector experiences a strong suppression of PBH formation due to the
elevated collapse threshold and reduced horizon mass. This sectoral asymmetry is a distinctive
and testable prediction of RDCC and forms the basis for the modified collapse threshold
developed in Section~4.

% ========================================================
% 4 MODIFIED COLLAPSE THRESHOLD IN RDCC
% ========================================================

\section{Modified Collapse Threshold in RDCC}

Primordial black holes form when large density fluctuations re-enter the horizon during
radiation domination and exceed a critical collapse threshold. In standard cosmology, this
threshold is of order
\begin{equation}
\delta_{\mathrm{std}} \sim 0.4 - 0.5,
\end{equation}
and PBH formation requires either enhanced primordial fluctuations or non-Gaussian features.

In the Relaxation-Driven Cyclic Cosmology (RDCC), the early-Universe dynamics are modified
by the universal relaxation mechanism introduced in Companion~X. The relaxation rate
\(\Gamma(a) = \alpha(a) H(a)\) alters the growth of super-horizon perturbations and increases the
effective collapse threshold in the visible sector. This modification arises from the additional
damping term in the perturbation equation of motion.

% --------------------------------------------------------
% 4.1 Relaxation-Induced Damping
% --------------------------------------------------------

\subsection{Relaxation-Induced Damping}

The evolution of super-horizon perturbations in RDCC follows the modified equation
\begin{equation}
\ddot{\delta} + (2H + \Gamma)\dot{\delta} - 2H^{2}\delta = 0,
\end{equation}
where the additional damping term \(\Gamma(a)\) arises from the inter-sector relaxation dynamics.
Using the relation
\begin{equation}
\Gamma(a) = \alpha(a) H(a),
\end{equation}
the equation becomes
\begin{equation}
\ddot{\delta} + (2 + \alpha)H\dot{\delta} - 2H^{2}\delta = 0.
\end{equation}

In the relaxation-dominated regime, \(\Gamma \sim \alpha H\), the growth of \(\delta\) is suppressed by a
factor that can be parametrized as an effective enhancement of the collapse threshold. Solving
the equation yields the scaling
\begin{equation}
\delta(a) \propto a^{-\alpha/2},
\end{equation}
implying that the initial amplitude must be larger by a factor \(a^{\alpha/2}\) to reach the standard
collapse threshold at horizon re-entry.

% --------------------------------------------------------
% 4.2 Relaxation-Corrected Collapse Threshold
% --------------------------------------------------------

\subsection{Relaxation-Corrected Collapse Threshold}

This motivates the parametrization
\begin{equation}
\delta_{c}^{(A/B)}(a) = \delta_{\mathrm{std}}
\left[ 1 + \beta\, \alpha_{A/B}(a) \right],
\end{equation}
where \(\beta = \mathcal{O}(1)\) captures the dependence on the perturbation profile. Because the
CPT-conjugate sector is colder, the relaxation coefficient satisfies
\begin{equation}
\alpha_{B}(a) < \alpha_{A}(a),
\end{equation}
and therefore
\begin{equation}
\delta_{c}^{(B)} < \delta_{c}^{(A)}.
\end{equation}

% --------------------------------------------------------
% 4.3 Exponential Sensitivity of PBH Formation
% --------------------------------------------------------

\subsection{Exponential Sensitivity of PBH Formation}

PBH formation is exponentially sensitive to the collapse threshold. Writing the PBH mass
fraction as
\begin{equation}
\beta_{\mathrm{PBH}}^{(A/B)}(M)
\propto
\exp\!\left[
-\frac{\delta_{c}^{2\,(A/B)}}{2\sigma^{2}(M)}
\right],
\end{equation}
makes clear that even a small relative change
\begin{equation}
\frac{\Delta \delta_{c}}{\delta_{c}}
\sim
\beta\, \alpha_{A/B}(a),
\end{equation}
induces an exponentially large difference in \(\beta_{\mathrm{PBH}}\).

This is the origin of the strong sectoral PBH asymmetry in RDCC:
\begin{equation}
\beta_{\mathrm{PBH}}^{(B)} \gg \beta_{\mathrm{PBH}}^{(A)}.
\end{equation}

The visible sector is exponentially suppressed, while the CPT sector naturally forms PBHs.
This sectoral asymmetry is a direct and testable prediction of the RDCC framework and forms
the basis for the PBH mass spectrum developed in Section~5.

% ========================================================
% 5 THERMODYNAMICS AND DYNAMICS OF CPT-SECTOR PBHs
% ========================================================

\section{Thermodynamics and Dynamics of CPT-Sector PBHs}

PBHs formed in the CPT-conjugate sector inherit the thermodynamic and dynamical properties
of that sector. Because the CPT sector evolves with a lower temperature, reduced relaxation
coefficient, and modified perturbation dynamics, the resulting PBHs differ systematically from
those that would form in the visible sector. Three features are particularly important.

% --------------------------------------------------------
% 5.1 Temperature Dependence
% --------------------------------------------------------

\subsection{Temperature Dependence}

The colder temperature of sector~B implies:

\begin{itemize}
\item a slower expansion rate,
\item a larger horizon mass at a given scale factor,
\item reduced pressure support,
\item a lower collapse threshold.
\end{itemize}

These effects combine to enhance PBH formation relative to the visible sector. The larger
horizon mass directly shifts the PBH mass spectrum toward higher masses, while the reduced
pressure support lowers the threshold for collapse.

% --------------------------------------------------------
% 5.2 Relaxation History
% --------------------------------------------------------

\subsection{Relaxation History}

The relaxation coefficient satisfies
\begin{equation}
\alpha_{B}(a) < \alpha_{A}(a),
\end{equation}
implying reduced damping of super-horizon perturbations in sector~B. This allows overdensities
to grow more efficiently prior to horizon re-entry. The relaxation mechanism therefore plays a
dual role:

\begin{enumerate}[label=(\alph*)]
\item it suppresses PBH formation in the visible sector by increasing the effective collapse threshold,
\item it enhances PBH formation in the CPT sector by reducing damping and lowering the threshold.
\end{enumerate}

This sectoral asymmetry is a direct consequence of the bipartite global state and the universal
relaxation dynamics.

% --------------------------------------------------------
% 5.3 Mass Spectrum and Abundance
% --------------------------------------------------------

\subsection{Mass Spectrum and Abundance}

The PBH mass spectrum in sector~B is shifted toward higher masses due to the larger horizon
mass:
\begin{equation}
M_{\mathrm{PBH}}^{(B)} \simeq 4\, M_{\mathrm{PBH}}^{(A)}.
\end{equation}

The abundance is exponentially enhanced:
\begin{equation}
\beta_{\mathrm{PBH}}^{(B)}
\propto
\exp\!\left[
-\frac{\delta_{c}^{2\,(B)}}{2\sigma^{2}}
\right],
\qquad
\delta_{c}^{(B)} < \delta_{c}^{(A)}.
\end{equation}

Together, these features establish PBHs as natural CPT-sector objects in RDCC, arising from
the interplay between relaxation dynamics, sectoral thermodynamics, and the bipartite global
state.

% ========================================================
% 6 OBSERVATIONAL SIGNATURES
% ========================================================

\section{Observational Signatures}

The sectoral PBH asymmetry predicted by the Relaxation-Driven Cyclic Cosmology (RDCC)
leads to a number of distinctive observational signatures. These signatures arise from the
suppression of PBHs in the visible sector, the enhanced PBH abundance in the CPT-conjugate
sector, and the shift of the PBH mass spectrum toward higher masses in sector~B. Together,
they provide a powerful observational handle on the relaxation-driven dynamics that define the
RDCC framework.

% --------------------------------------------------------
% 6.1 Suppressed PBH Abundance in the Visible Sector
% --------------------------------------------------------

\subsection{Suppressed PBH Abundance in the Visible Sector}

The relaxation-induced enhancement of the collapse threshold in sector~A,
\begin{equation}
\delta_{c}^{(A)} = \delta_{\mathrm{std}} \left[ 1 + \beta\, \alpha_{A}(a) \right],
\end{equation}
leads to an exponential suppression of PBH formation. This prediction is consistent with
microlensing, CMB, and SGWB constraints, all of which show no evidence for a significant PBH
population in the visible sector.

% --------------------------------------------------------
% 6.2 Enhanced PBH Abundance in the CPT Sector
% --------------------------------------------------------

\subsection{Enhanced PBH Abundance in the CPT Sector}

Because \(\alpha_{B}(a) < \alpha_{A}(a)\), the collapse threshold in sector~B is lower,
\begin{equation}
\delta_{c}^{(B)} < \delta_{c}^{(A)},
\end{equation}
leading to an exponentially larger PBH abundance. These PBHs are not directly observable, but
their existence is a robust prediction of the bipartite global state and the relaxation dynamics.

% --------------------------------------------------------
% 6.3 Sectoral Mass Asymmetry
% --------------------------------------------------------

\subsection{Sectoral Mass Asymmetry}

The colder temperature of sector~B implies a larger horizon mass at collapse,
\begin{equation}
M_{\mathrm{PBH}}^{(B)} \simeq 4\, M_{\mathrm{PBH}}^{(A)},
\end{equation}
leading to heavier PBHs in the CPT sector. This mass asymmetry is a direct consequence of the
sectoral temperature ratio derived in Companion~X.

% --------------------------------------------------------
% 6.4 Connection to Tensor Modulation
% --------------------------------------------------------

\subsection{Connection to Tensor Modulation}

The sectoral PBH asymmetry is part of a broader pattern of inter-sectoral imprints in RDCC.
Companion~II shows that the same bipartite structure and relaxation dynamics induce a
modulation of the tensor sector, affecting the primordial gravitational-wave background. Both
effects—the CPT-sector PBH population and the tensor modulation—arise from the same
underlying inter-sectoral coupling and the temperature asymmetry between the sectors. PBHs
therefore provide a complementary probe of the physics that shapes the tensor sector in RDCC.

% --------------------------------------------------------
% 6.5 Gravitational-Wave Signatures
% --------------------------------------------------------

\subsection{Gravitational-Wave Signatures}

PBHs contribute to the stochastic gravitational-wave background (SGWB) through mergers and
collapse. The suppression in sector~A implies:
\begin{itemize}
\item a reduced merger rate,
\item no PBH-induced SGWB in the visible sector,
\item possible indirect signatures from sector~B.
\end{itemize}

These signatures complement the relaxation-induced tensor suppression developed in
Companion~XV.

% --------------------------------------------------------
% 6.6 Microlensing Constraints
% --------------------------------------------------------

\subsection{Microlensing Constraints}

The strong suppression of PBHs in sector~A implies that microlensing surveys should observe no
significant PBH population—consistent with current data from EROS, OGLE, Subaru/HSC, and
Roman forecasts.

% --------------------------------------------------------
% 6.7 Modified PBH Mass Spectrum
% --------------------------------------------------------

\subsection{Modified PBH Mass Spectrum}

The relaxation-corrected mass function,
\begin{equation}
\beta_{\mathrm{PBH}}^{(A)}(M)
\propto
\exp\!\left[
-\frac{\delta_{c}^{2\,(A)}}{2\sigma^{2}(M)}
\right],
\end{equation}
predicts a suppressed and shifted mass spectrum in sector~A. The visible sector therefore
exhibits neither a significant PBH abundance nor a broad mass distribution.

% --------------------------------------------------------
% 6.8 Consistency with RDCC Temperature Asymmetry
% --------------------------------------------------------

\subsection{Consistency with RDCC Temperature Asymmetry}

The sectoral PBH asymmetry provides an independent test of the temperature ratio
\begin{equation}
\frac{T_{B}}{T_{A}} \simeq 0.5,
\end{equation}
derived in Companion~X. The PBH sector therefore offers a complementary observational probe
of the same relaxation-driven thermodynamics that shape the neutrino and dark matter sectors.

% ========================================================
% 7 FORECASTS
% ========================================================

\section{Forecasts}

The sectoral PBH asymmetry predicted by the Relaxation-Driven Cyclic Cosmology (RDCC)
provides multiple observational avenues for testing the framework. Future surveys will be able
to probe the PBH sector through gravitational waves, microlensing, and the matter power
spectrum. These forecasts highlight the complementarity between PBH physics, tensor
observables, and sectoral thermodynamics.

% --------------------------------------------------------
% 7.1 Gravitational-Wave Forecasts
% --------------------------------------------------------

\subsection{Gravitational-Wave Forecasts}

Space-based interferometers such as LISA, DECIGO, and BBO will be sensitive to the stochastic
gravitational-wave background (SGWB) generated by PBH mergers in the CPT sector. While
these PBHs are not directly observable, their merger history may leave detectable imprints in
the SGWB. The predicted signatures include:

\begin{itemize}
\item a sectoral enhancement of the high-frequency SGWB due to heavier PBHs in sector~B,
\item a merger-driven contribution consistent with the relaxation-induced tensor suppression
developed in Companion~XV,
\item a characteristic frequency dependence linked to the horizon mass in the CPT sector.
\end{itemize}

These signatures provide an indirect but powerful probe of the CPT-sector PBH population.

% --------------------------------------------------------
% 7.2 Microlensing Forecasts
% --------------------------------------------------------

\subsection{Microlensing Forecasts}

Microlensing surveys will continue to tighten constraints on PBHs in the visible sector. The
RDCC prediction of strong suppression in sector~A is consistent with current data and will be
further tested by future surveys such as Roman and LSST. The expected signatures are:

\begin{itemize}
\item no detectable PBH population in the visible sector,
\item increasingly stringent upper bounds on the PBH mass fraction in sector~A,
\item consistency with the relaxation-corrected collapse threshold derived in Section~4.
\end{itemize}

A null detection in microlensing is therefore a natural outcome of the RDCC framework.

% --------------------------------------------------------
% 7.3 Large-Scale-Structure Forecasts
% --------------------------------------------------------

\subsection{Large-Scale-Structure Forecasts}

The absence of PBH-induced anomalies in the matter power spectrum provides an additional
test of the RDCC PBH sector. Future surveys such as Euclid and CMB-S4 will be able to probe:

\begin{itemize}
\item small-scale suppression consistent with the relaxation-modified perturbation dynamics,
\item the absence of PBH-induced isocurvature modes in the visible sector,
\item consistency between the PBH sector and the sectoral temperature ratio \(T_{B}/T_{A} \simeq 0.5\).
\end{itemize}

These forecasts provide a cross-check between PBH physics, tensor observables, and the
relaxation-driven thermodynamics that define the RDCC framework.

% ========================================================
% 8 CONCLUSION
% ========================================================

\section{Conclusion}

In this Companion we have developed the primordial black hole (PBH) sector of the
Relaxation-Driven Cyclic Cosmology (RDCC), introducing a sectoral interpretation in which
PBHs arise naturally and predominantly in the CPT-conjugate sector. The colder temperature,
reduced relaxation-induced damping, and modified perturbation dynamics in sector~B enhance
PBH formation, while the visible sector experiences a strong suppression due to the elevated
collapse threshold and reduced horizon mass.

These effects arise directly from the bipartite CPT-symmetric global state and the universal
relaxation mechanism introduced in Companion~X. The resulting sectoral PBH asymmetry is a
robust and testable prediction of RDCC, linking PBH physics to the same relaxation-driven
thermodynamics that shape the neutrino and dark matter sectors (Companion~XIII) and the
tensor sector (Companion~XV).

The PBH mass spectrum in the CPT sector is shifted toward higher masses, and the PBH
abundance is exponentially enhanced relative to the visible sector. These features provide a
distinctive observational target, with signatures in the stochastic gravitational-wave background,
microlensing constraints, and the matter power spectrum. Future surveys such as LISA, DECIGO,
BBO, Roman, LSST, Euclid, and CMB-S4 will be able to probe these predictions.

With the development of the PBH sector, RDCC now offers a coherent and internally
consistent description of early-Universe thermodynamics, perturbation dynamics, and sectoral
structure. PBHs emerge as a natural and robust component of the CPT sector and a key probe
of relaxation-driven cosmology.


\end{document}
